% ---------------------------------------------------------------------------
%  FIGURE 2 -- Experimental pipeline: one trained detector, four protocols.
%
%  Rebuilt in TikZ from figures/source/fig2_pipeline_raw.jpeg. The pictograms
%  are the originals (scripts/extract_figure_icons.py); the heavy-bordered
%  panels, the progressive greys, the chevron stage chain and the hatched
%  inactive arm follow that render. Every word is set in LaTeX.
%
%  Corrections against the render:
%   * it carried three layout captions -- "LEFT SPINE", "RIGHT SIDE",
%     "RUNTIME / LATENCY ARM". Those were instructions to the renderer, not
%     labels, and they are removed along with the title.
%   * it labelled a stage "measured latency" and drew an offered-load versus
%     queueing-delay curve. No load sweep was performed and the latency is
%     emulated, so the runtime panel now reports only the stage decomposition
%     that was measured, and is marked EMULATED.
% ---------------------------------------------------------------------------
\providecommand{\ficontwo}[2]{%
  \raisebox{-0.45\height}{\includegraphics[height=#2]{../figures/icons/#1.png}}}

\definecolor{pInk}{HTML}{1F242B}
\definecolor{pG1}{HTML}{F4F6F8}\definecolor{pG2}{HTML}{E4E9EE}
\definecolor{pG3}{HTML}{D2DAE2}\definecolor{pG4}{HTML}{BCC7D2}
\definecolor{pEdge}{HTML}{6B7785}

\begin{tikzpicture}[
  font=\scriptsize,
  every node/.style={inner sep=1pt, outer sep=0pt},
  spine/.style ={draw=pInk, fill=white, rounded corners=2pt, line width=0.9pt,
                 align=left, inner sep=4pt, text width=3.05cm,
                 minimum height=1.05cm},
  % rung fill and border both deepen as the protocol nears deployment
  rungA/.style ={draw=pEdge, line width=0.8pt, fill=pG1, rounded corners=2pt,
                 align=left, inner sep=4.5pt, text width=5.55cm,
                 minimum height=1.12cm},
  rungB/.style ={rungA, draw=pEdge, line width=1.05pt, fill=pG2},
  rungC/.style ={rungA, draw=pInk!80, line width=1.35pt, fill=pG3},
  rungD/.style ={rungA, draw=pInk, line width=1.7pt, fill=pG4},
  blkd/.style  ={draw=pEdge!85, dashed, rounded corners=2pt, line width=0.9pt,
                 pattern=north east lines, pattern color=pEdge!45,
                 align=center, inner sep=5pt, text width=5.55cm,
                 dash pattern=on 2.4pt off 1.8pt},
  chev/.style  ={draw=pEdge, fill=pG2, line width=0.8pt, align=center,
                 signal, signal to=east, signal pointer angle=120,
                 inner xsep=3.5pt, inner ysep=3pt, font=\tiny},
  chevhot/.style={chev, draw=pInk, fill=pG4, line width=1.4pt},
  arr/.style   ={-{Latex[length=1.9mm]}, draw=pInk!75, line width=0.75pt},
  arrx/.style  ={-{Latex[length=1.9mm]}, draw=pEdge!70, dashed,
                 line width=0.7pt, dash pattern=on 2pt off 1.6pt},
  num/.style   ={circle, draw=pInk, fill=white, inner sep=0.8pt,
                 font=\scriptsize\bfseries, minimum size=4.2mm, line width=0.7pt},
  cost/.style  ={font=\scriptsize\bfseries, text=pInk, align=left, anchor=west},
  panel/.style ={draw=pInk, fill=white, rounded corners=2pt, line width=0.9pt,
                 inner sep=5pt, align=center, text width=4.25cm},
]

% ===== training spine ======================================================
\node[spine] (src) at (0,7.35)
  {\textbf{$D_A$: O-RAN testbed} \\ {\tiny Zeek flow records $+$ radio KPIs}};
\node at (2.18,7.35) {\ficontwo{src_testbed}{4.6mm}};

\node[spine] (fsp) at (0,5.55)
  {\textbf{Shared feature space} \\ {\tiny 15 concepts / 18 columns} \\
   {\tiny semantic harmonisation}};
\node at (2.22,5.55) {\ficontwo{schema}{5.4mm}};

\node[spine] (spl) at (0,3.75)
  {\textbf{Group-disjoint split} \\ {\tiny \texttt{src\_ip}, 20 split seeds}};
\node at (2.24,3.75) {\ficontwo{grouplock}{4.8mm}};

\node[spine, minimum height=1.62cm] (trn) at (0,1.62)
  {\textbf{Train} \\ {\tiny 8 architectures, including 2 trivial floors} \\
   {\tiny L1\;\; majority, stratified} \\
   {\tiny L2\;\; logreg, tree, RF, XGB, HGB} \\
   {\tiny L3\;\; MLP}};

\foreach \a/\b in {src/fsp, fsp/spl, spl/trn} {\draw[arr] (\a) -- (\b);}

% ===== one trained detector feeds every rung ===============================
\draw[arr, line width=1.5pt] (trn.east) -- (3.42,1.62);
\node[font=\tiny, align=center, text=pInk] at (3.44,8.02)
  {ONE TRAINED DETECTOR, \\ FOUR PROTOCOLS};
\draw[pInk!75, line width=0.75pt] (3.42,1.62) -- (3.42,7.35);

% ===== the evaluation ladder ===============================================
\node[rungA] (r1) at (7.55,7.35)
  {\textbf{Held-out $D_A$}, group-disjoint \\
   {\tiny in-distribution reference: the number this literature usually reports}};
\node[rungB] (r2) at (7.55,5.55)
  {\textbf{Time-disjoint $D_A$} \\
   {\tiny train on the earliest capture runs, test on the latest; 52.96 days of
   real elapsed time}};
\node[rungC] (r3) at (7.55,3.75)
  {\textbf{$D_B$: 5G-NIDD} \\
   {\tiny an independently collected deployment, evaluated once; every access
   logged}};
\node[rungD] (r4) at (7.55,1.62)
  {\textbf{Operational analysis} \\
   {\tiny base rate $\pi$, decision threshold $\tau$, benign arrival rate
   $\lambda_b$, alert burden}};

\foreach \n/\y in {1/7.35, 2/5.55, 3/3.75, 4/1.62}
  {\node[num] at (3.98,\y) {\n};
   \draw[arr] (4.24,\y) -- (4.62,\y);}

\node[cost] at (10.62,7.35) {reference};
\node[cost] at (10.62,5.55) {$-0.27\ F_{1}$};
\node[cost, text width=1.95cm] at (10.62,3.75) {to the trivial floor};
\node[cost, text width=2.45cm] at (10.62,1.62) {PPV $0.93 \!\to\! 0.006$};
\node at (10.18,2.72) {\ficontwo{distrib}{4.4mm}};

% ===== the arm that was not executed =======================================
\node[blkd, minimum height=0.82cm] (blk) at (7.55,-0.42)
  {\textbf{LIVE NEAR-RT RIC --- NOT EXECUTED}\\[1pt]{\tiny see Limitations}};
\draw[arrx] (3.42,1.62) -- (3.42,-0.42) -- (blk.west);

% ===== runtime decomposition, emulated =====================================
\node[panel] (rt) at (14.95,7.28)
  {\textbf{RUNTIME DECOMPOSITION} \\[1pt]
   {\tiny\textsc{emulated} --- no RIC in the path}};

\begin{scope}[shift={(12.80,5.86)}]
  \node[chev]    (s0) at (0.42,0) {E2 \\ indic.};
  \node[chevhot] (s1) at (1.80,0) {feature \\ constr.};
  \node[chev]    (s2) at (3.10,0) {model \\ infer.};
  \node[chev]    (s3) at (4.26,0) {control \\ encode};
\end{scope}

\node[align=center, font=\tiny, text=pInk, text width=4.3cm] at (14.95,4.52)
  {$t_{\mathrm{ind}} + t_{\mathrm{feat}} + t_{\mathrm{inf}} + t_{\mathrm{act}}$ \\[2pt]
   \textbf{FEATURE CONSTRUCTION DOMINATES} \\[1pt]
   0.058\,ms per flow, against 1--25\,$\mu$s of inference in a deployed
   implementation};

\node[panel, text width=4.25cm] at (14.95,2.28)
  {\textbf{HOW TO READ THIS FIGURE} \\[2pt]
   {\tiny Every step down the ladder is closer to what a deployment
   experiences, and every step costs more than the choice of architecture
   does. Rung~1 alone does not establish deployment behaviour.}};

\end{tikzpicture}
